\section{Coronel Milena Mejía --- Defensa Civil}

\subsection{Información básica}
Coronel Milena Mejía, directora de la Defensa Civil. Cuenta con experiencia en la primera respuesta a un terremoto desde Quibdó y Cali, y también estuvo en Montería.

\subsection{Anotaciones}
\begin{itemize}
  \item Ser cordial, profesional y respetuoso en todo momento.
  \item Presentarse de parte del Coronel Ubaque, quien ya habló con ella.
  \item Agradecer de entrada el tiempo que va a dedicar a apoyar el proyecto.
  \item Explicar de forma breve que la entrevista tiene un propósito académico.
  \item Cuidado con las preguntas.
\end{itemize}

\subsection{Perfil profesional}
\begin{enumerate}
  \pregunta ¿Cuál ha sido su trayectoria en la dirección y coordinación de operaciones de emergencia?
  \pregunta ¿Con qué entidades y equipos debe articularse normalmente durante una emergencia?
  \pregunta Desde su experiencia, ¿cuáles son las capacidades más importantes que debe tener un equipo de gestión de emergencias?
\end{enumerate}

\subsection{Perspectiva operativa y tecnológica}
\begin{enumerate}
  \pregunta Durante las primeras horas de una emergencia, ¿qué información necesita con mayor urgencia y qué dificultades encuentra al anticipar la evolución de un incendio o al definir prioridades de respuesta?
  \pregunta ¿Qué herramientas, mapas, modelos o sistemas de información usan actualmente para apoyar la coordinación?
\end{enumerate}


% ---------------------------------------------------------
% CONTENIDO AÑADIDO
% ---------------------------------------------------------

\subsection{Flujo actual de respuesta y toma de decisiones}

\begin{enumerate}
  \pregunta Cuando se reporta inicialmente una emergencia, ¿cómo es el proceso desde que se recibe la alerta hasta las primeras decisiones operativas, quiénes participan y cómo se distribuyen las responsabilidades?


  \pregunta ¿De qué fuentes reciben la información inicial sobre el evento y qué información suele estar disponible inmediatamente o tardar más tiempo en obtenerse?




\end{enumerate}


\subsection{Información necesaria durante incendios forestales}

\begin{enumerate}


  \pregunta ¿Cómo identifican actualmente las zonas, comunidades o infraestructuras que podrían resultar amenazadas si el incendio continúa avanzando?



  \pregunta ¿Recuerda alguna situación en la que disponer de una mejor estimación de la evolución de una emergencia hubiera permitido tomar una decisión diferente?
\end{enumerate}


\subsection{Limitaciones de los procesos actuales}

\begin{enumerate}
  \pregunta ¿Cuál considera que es actualmente la principal dificultad para obtener una visión completa de lo que está ocurriendo durante una emergencia?

  \pregunta ¿Existen problemas relacionados con la actualización o precisión de los datos, o dificultades para compartir información entre las distintas entidades involucradas?


  \pregunta ¿La conectividad en las zonas afectadas suele limitar el uso de herramientas tecnológicas?


  \pregunta ¿Qué procesos considera que consumen más tiempo del necesario durante la planeación o coordinación de una emergencia?
\end{enumerate}


\subsection{Validación de una posible herramienta}

Antes de realizar las siguientes preguntas se explicará brevemente que el proyecto
estudia la posibilidad de utilizar escenarios de incendios previamente simulados
para comparar una emergencia real con situaciones similares y ofrecer información
rápidamente como apoyo a la toma de decisiones.

\begin{enumerate}
  \pregunta Pensando en una herramienta de este tipo, ¿en qué momento de una emergencia y para qué decisiones concretas considera que podría aportar mayor valor?

  \pregunta ¿Qué información debería mostrar obligatoriamente para resultar útil?

  \pregunta ¿Qué información considera que podría resultar innecesaria o dificultar la interpretación?





\end{enumerate}


\subsection{Cierre}

\begin{enumerate}


  \pregunta ¿Qué aspecto de una solución como la planteada considera que deberíamos investigar con mayor profundidad antes de desarrollarla?

  \pregunta ¿Qué otras personas, entidades o perfiles profesionales considera que deberíamos entrevistar para comprender mejor este problema?
\end{enumerate}



% =========================================================
% ENTREVISTA 2
% =========================================================

\clearpage
